% mooie C++
% http://archive.cs.uu.nl/mirror/CTAN/macros/latex/contrib/relsize/relsize-doc.pdf
\usepackage{relsize}
\iftoggle{charter}{
    \newcommand{\Cpp}{C\nolinebreak\raisebox{.2ex}[0pt][0pt]{\smaller{+\nolinebreak+}}}
}{
\newcommand{\Cpp}{C\nolinebreak\raisebox{.3ex}[0pt][0pt]{\smaller{+\nolinebreak+}}}
}

% define kleuren die in listings gebruikt worden
\definecolor{lstlightgray}{RGB}{128,128,128}
\definecolor{lstlightergray}{RGB}{169,169,199}
\definecolor{lstdarkred}{RGB}{191,0,0}
\definecolor{lstdarkerred}{RGB}{128,0,0}
\definecolor{lstpurple}{RGB}{127,0,116}
\definecolor{lstmatlabpurple}{RGB}{159,31,242}
\definecolor{lstdarkgreen}{RGB}{0,128,0}
\definecolor{lstdarkergreen}{RGB}{0,100,0}
\definecolor{lstturquoisegreen}{RGB}{63,127,95}
\definecolor{lstmatlabturquoisegreen}{RGB}{33,138,77}
\definecolor{lstblue}{RGB}{42,0,255}
\definecolor{lstblueer}{RGB}{0,0,255}
\definecolor{llightblue}{RGB}{0,149,255}
\definecolor{lstdarkblue}{RGB}{0,0,128}
\definecolor{lstdarkerblue}{RGB}{15,15,127}

% http://archive.cs.uu.nl/mirror/CTAN/macros/latex/contrib/listings/listings.pdf
\usepackage{listings}
\lstset{
    literate=
        {ä}{{\"a}}1 
        {ë}{{\"e}}1 
        {ï}{{\"i}}1 
        {ö}{{\"o}}1
        {ü}{{\"u}}1
        {á}{{\'a}}1
        {é}{{\'e}}1
        {í}{{\'i}}1
        {ó}{{\'o}}1
        {ú}{{\'u}}1
        {à}{{\`a}}1
        {è}{{\`e}}1
        {ì}{{\`i}}1
        {ò}{{\`o}}1
        {ù}{{\`u}}1
        {â}{{\^a}}1
        {ê}{{\^e}}1
        {î}{{\^i}}1
        {ô}{{\^o}}1
        {û}{{\^u}}1
        {©}{{\copyright}}1,
    basicstyle=\linespread{1.1}\ttfamily\small,
    tabsize=4,
    showstringspaces=false,
    upquote=true,
    numbers=none,
    aboveskip=.1\baselineskip,
    belowskip=0pt,
    abovecaptionskip=.3\baselineskip,
    belowcaptionskip=-.2\baselineskip,
    includerangemarker=false,
    prebreak={\raisebox{0ex}[0ex][0ex]{\textcolor{black}{\ensuremath{\hookleftarrow}}}}, 
    postbreak={\raisebox{0ex}[0ex][0ex]{\textcolor{black}{\ensuremath{\space\hookrightarrow\space}}}}, 
    breaklines=true, 
    breakatwhitespace=true, 
    breakindent=0.1em,
    numberbychapter=true,
    captionpos=b,
    keepspaces=true,
}

\lstdefinestyle{linenumbers}{
    numbers=left, 
    stepnumber=1, 
    numberstyle=\tiny\textcolor{hrred}, 
    numbersep=1em
}

% C 
% corlor scheme from Code Composer Studio
\newcommand{\lstcstring}{\color{lstblue}}
\newcommand{\lstckeyword}{\color{lstpurple}\bfseries}
\newcommand{\lstcincludefile}{\color{lstblue}}
\newcommand{\lstccomment}{\color{lstturquoisegreen}}
\lstdefinestyle{cstyle}{
    language=C,
    morekeywords={__interrupt, __asm, interrupt, asm, __attribute__,true,false},
    morekeywords=[2]{<assert.h>,<complex.h>,<ctype.h>,<errno.h>,<fenv.h>,<float.h>,<inttypes.h>,<iso646.h>,<limits.h>,<locale.h>,<math.h>,<setjmp.h>,<signal.h>,<stdalign.h>,<stdarg.h>,<stdatomic.h>,<stdbool.h>,<stddef.h>,<stdint.h>,<stdio.h>,<stdlib.h>,<stdnoreturn.h>,<string.h>,<tgmath.h>,<threads.h>,<time.h>,<uchar.h>,<wchar.h>,<wctype.h>},
    stringstyle=\lstcstring,
    keywordstyle=\lstckeyword,
    alsoletter={<>.},
	keywordstyle=[2]\lstcincludefile,
    identifierstyle=,
    commentstyle=\lstccomment,
}

% C++
% corlor scheme from Code Composer Studio
\newcommand{\lstcppstring}{\color{lstblue}}
\newcommand{\lstcppkeyword}{\color{lstpurple}\bfseries}
\newcommand{\lstcppincludefile}{\color{lstblue}}
\newcommand{\lstcppcomment}{\color{lstturquoisegreen}}
\lstdefinestyle{cppstyle}{
    language=C++,
    morekeywords={alignas,alignof,char16_t,char32_t,concept,constexpr,decltype,noexcept,nullptr,requires,static_assert,thread_local,override,final},
    morekeywords=[2]{<cstdlib>,<csignal>,<csetjmp>,<cstdarg>,<typeinfo>,<typeindex>,<type_traits>,<bitset>,<functional>,<utility>,<ctime>,<chrono>,<cstddef>,<initializer_list>,<tuple>,<any>,<optional>,<variant>,<compare>,<version>,<new>,<memory>,<scoped_allocator>,<memory_resource>,<climits>,<cfloat>,<cstdint>,<cinttypes>,<limits>,<exception>,<stdexcept>,<cassert>,<system_error>,<cerrno>,<contract>,<cctype>,<cwctype>,<cstring>,<cwchar>,<cuchar>,<string>,<string_view>,<charconv>,<array>,<vector>,<deque>,<list>,<forward_list>,<set>,<map>,<unordered_set>,<unordered_map>,<stack>,<queue>,<span>,<iterator>,<ranges>,<algorithm>,<execution>,<cmath>,<complex>,<valarray>,<random>,<numeric>,<ratio>,<cfenv>,<bit>,<iosfwd>,<ios>,<istream>,<ostream>,<iostream>,<fstream>,<sstream>,<syncstream>,<strstream>,<iomanip>,<streambuf>,<cstdio>,<locale>,<clocale>,<codecvt>,<regex>,<atomic>,<thread>,<mutex>,<shared_mutex>,<future>,<condition_variable>,<filesystem>},
    stringstyle=\lstcppstring,
    keywordstyle=\lstcppkeyword,
    alsoletter={<>.},
	keywordstyle=[2]\lstcppincludefile,
    identifierstyle=,
    commentstyle=\lstcppcomment,
}

% LaTeX
% color scheme from TeXstudio
\newcommand{\lstlcs}{\color{lstdarkerred}}
\newcommand{\lstlbeginend}{\color{llightblue}}
\newcommand{\lstlenvironment}{\color{lstdarkblue}}
\newcommand{\lstlcomment}{\color{lstlightgray}}
\newcommand{\lstlmath}{\color{lstdarkgreen}}
\newcommand{\lstllabel}{\color{lstdarkgreen}}
\lstdefinestyle{lstyle}{
	language={[LaTeX]TeX},
	moretexcs={%
		% Alle latexcommando's is onbegonnen werk. Alle commando's die in style_NL_nl zijn gebruikt zijn hier opgenomen.
        iftoggle,cref,Cref,crefname,Crefname,cpageref,Cpageref,href,url,frontmatter,mainmatter,backmatter,baslineskip,noclub,nowidow,includegraphics,texorpdfstring,usetikzlibrary,tikz,draw,path,fill,filldraw,cllip,shade,shadedraw,foreach,node,coordinate,graph,matrix,calender,addplot,enquote,lstinputlisting,lstset,lstinline,arrayrulecolor,tabularxcolumn,textcite,si,SI,num,arraybackslash,
		% style_nl_NL
        myClass,modulecode,cursuscode,huidigeversie,titel,auteurs,documenttype,sectionfont,headerfont,footerfont,captionfont,titelpagina,versiehistorie,mailafko,figuur,hrfigure,tikzfiguur,hrtikzfigure,tikzinlinefiguur,hrtikzinlinefigure,schakeling,tabel,tabelmaxpaginabreedgeschaald,langetabel,tabelX,langetabelX,tabelXm,langetabelXm,USE,hyp,csecref,Csecref,pdfpagref,pdfdeelref,conclusie,handtekening,citepdfpag,citepdfdeel,dquote,squote,ebookfalse,ebooktrue,chartertrue,opensanstrue,timestrue,progurl,
		% styleCode
		Cpp,floatlstinput,lstinput,lstinputfragment,code,ccode,cppcode,lcode,pycode,mcode,acode,amcode,legcode,pinkycode,hcode,cinput,proglink	
	},
	texcsstyle=\lstlcs,
	deletetexcs={begin,end},
	moretexcs=[2]{begin,end},
	texcsstyle=[2]\lstlbeginend,
	morekeywords=[3]{%
        % standard environments
		array,center,description,document,enumerate,eqnarray,equation,figure,flushleft,flushright,itemize,list,minipage,picture,quotation,quote,tabbing,table,tabular,thebiblipgraphy,theorem,titlepage,verbatim,verse,
		%longtabu, tabu, // deprecated
		% Alle latexenvironments is onbegonnen werk. Alle environments die in style_NL_nl zijn gebruikt zijn hier opgenomen.
        tikzpicture,circuittikz,scope,axis,mdframed,tabularx,
        % style_nl_NL
		%itemizeintabu, // deprecated
		itemizeintabel,myFigure,myTable,opdracht,
		% styleCode
		floatlst,lst,lstfragment,clstShortInline,cpplstShortInline,llstShortInline,pylstShortInline,mlstShortInline,alstShortInline,amlstShortInline,leglstShortInline,pinkylstShortInline,coutput
	},
	keywordstyle=[3]\lstlenvironment,
	commentstyle=\lstlcomment,
	moredelim=[s][\lstlmath]{$}{$},	
	moredelim=[s][\lstlmath]{\\\[}{\\\]},	
}

% Python
% color scheme from Thonny http://thonny.org/
\newcommand{\lstpstring}{\color{lstdarkergreen}}
\newcommand{\lstpkeyword}{\color{lstpurple}\bfseries}
\newcommand{\lstpcomment}{\color{lstlightergray}}
\lstdefinestyle{pystyle}{
	language=Python,
	% redefined keywords for Python 3
	keywords={False,class,finally,is,return,None,continue,for,lambda,try,True,def,from,
		nonlocal,while,and,del,global,not,with,as,elif,if,or,yield,assert,else,import,
		pass,break,except,in,raise},
	keywords=[2]{},
	morekeywords={},
	stringstyle=\lstpstring,
	keywordstyle=\lstpkeyword,
	identifierstyle={},
	commentstyle=\lstpcomment,
}

% MATLAB
\newcommand{\lstmstring}{\color{lstmatlabpurple}}
\newcommand{\lstmkeyword}{\color{lstblueer}}
\newcommand{\lstmcomment}{\color{lstmatlabturquoisegreen}}
\lstdefinestyle{mstyle}{
    language=Matlab,
    morekeywords={},
    stringstyle=\lstmstring,
    keywordstyle=\lstmkeyword,
    identifierstyle=,
    commentstyle=\lstmcomment,
}

% AVR assembler
\lstdefinelanguage[AVR]{Assembler}{
    morekeywords=[1]{
        adc,add,adiw,andi,asr,bclr,bld,brbc,brbs,brcc,brcs,break,breq,brge,brhc,brhs,brid,brie,brlo,brlt,brmi,brne,brpl,brsh,brtc,brts,brvc,brvs,bset,bst,cbi,cbr,clc,clh,cli,cln,clr,cls,clt,clv,clz,com,count,cp,cpc,cpi,cpse,dec,eicall,eijmp,elpm,end,eor,espm,fmul,fmuls,fmulsu,in,inc,jmp,ld,ldd,ldi,lds,lsl,lsr,mov,movw,neg,nop,ori,out,pop,push,rcall,ret,reti,rjmp,rol,sbc,sbci,sbi,sbic,sbis,sbiw,sbr,sbrc,sbrs,sec,seh,sei,sen,ser,ses,sev,sez,sleep,st,std,sts,sub,subi,swap,tst,wdr
    },
    morekeywords=[2]{
        .byte,.cseg,.csegsize,.db,.device,.dseg,.dw,.endm,.equ,.eseg,.exit,.include,.includepath,.listmac,.macro,.nolist,.org,.set,.define,.undef,.ifdef,.ifndef,.else,.elseif,.elif,.message,.warning,.error,byte1,byte2,byte3,byte4,exp2,high,hwrd,log2,low,lwrd,page
    },
    alsoletter=.,
    sensitive=false,
    morecomment=[l]{\#},
    morecomment=[l]{;},
    morestring=[b]{"},
    morestring=[b]{'}
}[keywords,comments,strings]

\newcommand{\lstamnemonic}{\color{lstblueer}}
\newcommand{\lstacomment}{\color{lstdarkgreen}}
\lstdefinestyle{astyle}{
    language=[AVR]Assembler,
    stringstyle=,
    keywordstyle=[1]\lstamnemonic,
    keywordstyle=[2],
    identifierstyle=,
    commentstyle=\lstacomment,  
}

% MSP430 assembler
\lstdefinelanguage[MSP430]{Assembler}{
	morekeywords=[1]{
		adc,adc.b,adc.w,add,add.b,add.w,addc,addc.b,addc.w,and,and.b,and.w,bic,bic.b,bic.w,bis,bis.b,bis.w,bit,bit.b,bit.w,br,call,clr,clr.b,clr.w,clrc,clrn,clrz,cmp,cmp.b,cmp.w,dadc,dadc.b,dadc.w,dadd,dadd.b,dadd.w,dec,dec.b,dec.w,decd,decd.b,decd.w,dint,eint,inc,inc.b,inc.w,incd,incd.b,incd.w,inv,inv.b,inv.w,jc,jhs,jeq,jz,jge,jl,jmp,jn,jnc,jlo,jne,jnz,mov,mov.b,mov.w,nop,pop,pop.b,pop.w,push,push.b,push.w,ret,reti,rla,rla.b,rla.w,rlc,rlc.b,rlc.w,rra,rra.b,rra.w,rrc,rrc.b,rrc.w,sbc,sbc.b,sbc.w,setc,setn,setz,sub,sub.b,sub.w,subc,subc.b,subc.w,swpb,sxt,tst,tst.b,tst.w,xor,xor.b,xor.w,
	},
	morekeywords=[2]{
		.bss,.data,.sect,.text,.byte,.double,.float,.float,.long,.short,.word,.align,.align,.include,.common,.global,.if,.else,.endif,.set,.macro,.endm
	},
	alsoletter=.,
	sensitive=false,
	morecomment=[l]{;},
	morestring=[b]{"},
	morestring=[b]{'}
}[keywords,comments,strings]

\newcommand{\lstamstring}{\color{lstblue}}
\newcommand{\lstammnemonic}{\color{lstpurple}}
\newcommand{\lstamcomment}{\color{lstturquoisegreen}}
\lstdefinestyle{amstyle}{
	language=[MSP430]Assembler,
	stringstyle=\lstamstring,
	keywordstyle=[1]\lstammnemonic,
	keywordstyle=[2]\lstammnemonic,
	identifierstyle=,
	commentstyle=\lstamcomment,  
}

% LEG assembler
\lstdefinelanguage[LEG]{Assembler}{
	morekeywords=[1]{
		add,adds,addi,addis,and,andi,andis,ands,b,b.eq,b.ne,b.lt,b.le,b.gt,b.ge,b.lo,b.ls,b.hi,b.hs,bl,br,cbnz,cbz,eor,eori,ldur,ldurb,ldurh,ldursw,ldxr,lsl,lsr,movk,movz,orr,orri,stur,sturb,sturh,sturw,stxr,sub,subi,subis,subs,fadds,faddd,fcmps,fcmpd,fdivs,fdivd,fsubs,fsubd,ldurs,ldurd,mul,nop,sdiv,smulh,sturs,sturd,udiv,umulh,cmp,cmpi,lda,mov
	},
	morekeywords=[2]{
		x0,x1,x2,x3,x4,x5,x6,x7,x8,x9,x10,x11,x12,x13,x14,x15,x16,x17,x18,x19,x20,x21,x22,x23,x24,x25,x26,x27,x28,x29,x30,x31,xzr,sp,pc,fp,lr,ip0,ip1,w0,w1,w2,w3,w4,w5,w6,w7,w8,w9,w10,w11,w12,w13,w14,w15,w16,w17,w18,w19,w20,w21,w22,w23,w24,w25,w26,w27,w28,w29,w30,w31,wzr,d0,d1,d2,d3,d4,d5,d6,d7,d8,d9,d10,d11,d12,d13,d14,d15,d16,d17,d18,d19,d20,d21,d22,d23,d24,d25,d26,d27,d28,d29,d30,d31,s0,s1,s2,s3,s4,s5,s6,s7,s8,s9,s10,s11,s12,s13,s14,s15,s16,s17,s18,s19,s20,s21,s22,s23,s24,s25,s26,s27,s28,s29,s30,s31
	},
	morekeywords=[3]{
		.globl
	},
	morekeywords=[4]{
		\#0,\#1,\#2,\#3,\#4,\#5,\#6,\#7,\#8,\#9,\#10,\#11,\#12,\#13,\#14,\#15,\#16,\#17,\#18,\#19,\#20,\#21,\#22,\#23,\#24,\#25,\#26,\#27,\#28,\#29,\#30,\#31,\#32
	},
	alsoletter={:.\#},
	alsodigit={?},
	sensitive=false,
	morecomment=[l]{//},
	morestring=[b]{"},
	morestring=[b]{'}
	morecomment=[s]{/*}{*/},
}[keywords,comments,strings]

\newcommand{\lstlegstring}{\color{lstblue}}
\newcommand{\lstlegmnemonic}{\color{lstpurple}\bfseries}
%\newcommand{\lstlegregister}{\color{lstdarkerblue}\itshape}
% \itshape verwijderd omdat lettertype inconsolata geen cursief ondersteund
% zie: https://tex.stackexchange.com/questions/183204/inconsolata-italics
\newcommand{\lstlegregister}{\color{lstdarkerblue}}
\newcommand{\lstlegdirective}{\bfseries}
\newcommand{\lstlegimmediate}{\color{lstdarkred}}
\newcommand{\lstlegcomment}{\color{lstturquoisegreen}}
\lstdefinestyle{legstyle}{
	language=[LEG]Assembler,
    stringstyle=\lstlegstring,
	keywordstyle=\lstlegmnemonic,
	keywordstyle={[2]\lstlegregister},
	keywordstyle={[3]\lstlegdirective},
	keywordstyle={[4]\lstlegimmediate},
	identifierstyle=,
	commentstyle=\lstlegcomment,	
}

% Pinky assembler
\lstdefinelanguage[Pinky]{Assembler}{
	morekeywords=[1]{
		add.n,adds.n,adds.n,adds.n,ands.n,blx.n,bx.n,cbnz.n,cbz.n,cmp.n,cmp.n,eors.n,ldr.n,ldr.n,ldr.n,ldr.n,lsls.n,lsrs.n,movs.n,movs.n,nop.n,orrs.n,pop.n,push.n,str.n,str.n,str.n,sub.n,subs.n,subs.n,subs.n,beq.n,bne.n,bcs.n,blo.n,bcc.n,bhs.n,bmi.n,bpl.n,bvs.n,bvc.n,bhi.n,bls.n,bge.n,blt.n,bgt.n,ble.n,b.n,bal.n
	},
	morekeywords=[2]{
		r0,r1,r2,r3,r4,r5,r6,r7,r13,r14,r15,sp,msp,psp,lr,pc,lr,a1,a2,a3,a4,v1,v2,v3,v4,psr,primask,faultmask,baseptr,control
	},
	morekeywords=[3]{
		.globl,.global,.text,.data,.cpu,.thumb,.syntax,.thumb_func,.word
	},
	morekeywords=[4]{
		\#0,\#1,\#2,\#3,\#4,\#5,\#6,\#7,\#8,\#9,\#10,\#11,\#12,\#13,\#14,\#15,\#16,\#17,\#18,\#19,\#20,\#21,\#22,\#23,\#24,\#25,\#26,\#27,\#28,\#29,\#30,\#31,\#32
	},
	alsoletter={:.\#},
	alsodigit={?},
	sensitive=false,
	morecomment=[l]{//},
	morestring=[b]{"},
	morestring=[b]{'}
	morecomment=[s]{/*}{*/},
}[keywords,comments,strings]

\newcommand{\lstpinkystring}{\color{lstblue}}
\newcommand{\lstpinkymnemonic}{\color{lstpurple}\bfseries}
%\newcommand{\lstpinkyregister}{\color{lstdarkerblue}\itshape}
% \itshape verwijderd omdat lettertype inconsolata geen cursief ondersteund
% zie: https://tex.stackexchange.com/questions/183204/inconsolata-italics
\newcommand{\lstpinkyregister}{\color{lstdarkerblue}}
\newcommand{\lstpinkydirective}{\color{lstpurple}\bfseries}
\newcommand{\lstpinkyimmediate}{\color{lstdarkred}}
\newcommand{\lstpinkycomment}{\color{lstturquoisegreen}}
\lstdefinestyle{pinkystyle}{
	language=[Pinky]Assembler,
	stringstyle=\lstpinkystring,
	keywordstyle=\lstpinkymnemonic,
	keywordstyle={[2]\lstpinkyregister},
	keywordstyle={[3]\lstpinkydirective},
	keywordstyle={[4]\lstpinkyimmediate},
	identifierstyle=,
	commentstyle=\lstpinkycomment,	
}

\lstdefinestyle{nostyle}{
	stringstyle=,
	keywordstyle=,
	identifierstyle=,
	commentstyle=,
}

% Commando's \floatlstinput en \lstinput lezen code uit (een deel van) een bestand en plaatsen dit in een listing.
% #1 = (optional) extra key=value paren voor lstinputlisting b.v.
%      style = cstyle, cppstyle, pystyle, mstyle of astyle (C, C++, Python, MATLAB, AVR assembler)
%      style = linenumbers
%      linerange={1-6, 13-27}
% #2 = filenaam (moet in het directory ./progs staan)
% #3 = label=lst:#3
% #4 = titel
\newcommand{\floatlstinput}[4][]{%
\lstinputlisting[
    rangebeginprefix=//>,% Alleen nuttig voor C en C++ 
    rangeendprefix=//<,% Alleen nuttig voor C en C++ 
    float=!htbp, 
    label=lst:#3, 
    caption={#4},%    
    #1
]{progs/#2}%
}

% #1 = (optional) extra key=value paren voor lstinputlisting b.v.
%      style = cstyle, cppstyle, pystyle, mstyle of astyle (C, C++, Python, MATLAB, AVR assembler)
%      style = linenumbers
%      linerange={1-6, 13-27}
% #2 = filenaam (moet in het directory ./progs staan)
% #3 = label=lst:#3
% #4 = titel
\newcommand{\lstinput}[4][]{%
\lstinputlisting[
    aboveskip=.8\baselineskip,
    belowskip=.4\baselineskip,
    rangebeginprefix=//>,% Alleen nuttig voor C en C++ 
    rangeendprefix=//<,% Alleen nuttig voor C en C++ 
    label=lst:#3, 
    caption={#4},%    
    #1
]{progs/#2}%
}

% Commando \lstinputfragment leest code uit (een deel van) een bestand en plaatst dit in een codefragment.
% #1 = (optional) extra key=value paren voor lstinputlisting b.v.
%      style = cstyle, cppstyle, pystyle, mstyle of astyle (C, C++, Python, MATLAB, AVR assembler)
%      style = linenumbers
%      linerange={1-6, 13-27}
% #2 = filenaam (moet in het directory ./progs staan)
% geen titel, geen label
\newcommand{\lstinputfragment}[2][]{%
\lstinputlisting[
    aboveskip=.6\baselineskip,%
    rangebeginprefix=//>,% Alleen nuttig voor C en C++ 
    rangeendprefix=//<,% Alleen nuttig voor C en C++ 
    #1
]{progs/#2}%
}

% De code in de environments floatlst en lst wordt in een listing geplaatst.
% #1 = (optional) extra key=value paren voor lstset b.v.
%      style = cstyle, cppstyle, pystyle, mstyle of astyle (C, C++, Python, MATLAB, AVR assembler)
%      style = linenumbers
% #2 = label=lst:#2
% #3 = titel
\lstnewenvironment{floatlst}[3][]{%
    \lstset{
        float=!htbp, 
        label=lst:#2, 
        caption={#3},%   
        #1
    }%
}{%
}

% #1 = (optional) extra key=value paren voor lstset b.v.
%      style = cstyle, cppstyle, pystyle, mstyle of astyle (C, C++, Python, MATLAB, AVR assembler)
%      style = linenumbers
% #2 = label=lst:#2
% #3 = titel
\lstnewenvironment{lst}[3][]{%
    \lstset{
        aboveskip=.8\baselineskip,
        belowskip=.4\baselineskip,
        label=lst:#2, 
        caption={#3},%    
        #1
    }%
}{%
}

% De code in de environments lstfragment wordt in een codefragment geplaatst.
% #1 = (optional) extra key=value paren voor lstset b.v.
%      style = cstyle, cppstyle, pystyle, mstyle of astyle (C, C++, Python, MATLAB, AVR assembler)
%      style = linenumbers
% geen titel, geen label
\lstnewenvironment{lstfragment}[1][]{%
    \lstset{
        aboveskip=.6\baselineskip,
        #1
    }%
}{%
}

% \intextcodefont is nodig om \small te verwijderen
\newcommand{\intextcodefont}{\ttfamily}
\lstdefinestyle{lstinlinestyle}{
    basicstyle=\intextcodefont, 
    prebreak={}, 
    postbreak={}, 
    breakindent=0pt,
}

% Commando's \code, \lcode, \ccode, \cppcode, \pycode, \mcode en \acode kunnen gebruikt worden om inline code te definiëren.
% #1 = (optional) extra key=value paren voor lstinline.
% #2 = inline code
\newcommand{\code}[2][]{%
    \lstinline[
        style=lstinlinestyle,%
        #1
    ]@#2@%
}

\newcommand{\ccode}[2][]{%
    \lstinline[
        style=cstyle,
        style=lstinlinestyle,%
        #1
    ]@#2@%
}

\newcommand{\cppcode}[2][]{%
    \lstinline[
        style=cppstyle,
        style=lstinlinestyle,%
        #1
    ]@#2@%
}

\newcommand{\lcode}[2][]{%
	\lstinline[
		style=lstyle,
		style=lstinlinestyle,%
		#1
	]@#2@%
}

\newcommand{\pycode}[2][]{%
	\lstinline[
	style=pystyle,
	style=lstinlinestyle,%
	#1
	]@#2@%
}

\newcommand{\mcode}[2][]{%
    \lstinline[
        style=mstyle,
        style=lstinlinestyle,%
        #1
    ]@#2@%
}

\newcommand{\acode}[2][]{%
    \lstinline[
        style=astyle,
        style=lstinlinestyle,%
        #1
    ]@#2@%
}

\newcommand{\amcode}[2][]{%
	\lstinline[
	style=amstyle,
	style=lstinlinestyle,%
	#1
	]@#2@%
}

\newcommand{\legcode}[2][]{%
    \lstinline[
        style=legstyle,
        style=lstinlinestyle,%
        #1
    ]@#2@%
}

\newcommand{\pinkycode}[2][]{%
	\lstinline[
        style=pinkystyle,
        style=lstinlinestyle,%
	#1
	]@#2@%
}

% met \hcode kun je ervoor zorgen dat lange identifiers afgebroken kunnen worden. Mogelijke afbreekpunten met \\- markeren.
\newcommand{\hcode}[2][]{%
    \lstinline[
        style=lstinlinestyle,
        style=nostyle,
        literate={\\-}{}{0\discretionary{-}{}{}},%
        #1
    ]@#2@%
}

% handige afkorting |bla| voor \ccode{bla} maar pas op: werkt niet in footnotes, captions etc.
% werkt niet als je | in latex code gebruikt (b.v. in math).
% vandaar environment van gemaakt.
\newenvironment{clstShortInline}{%
    \lstMakeShortInline[
        style=cstyle,
        style=lstinlinestyle,
    ]{|}%
}{%
    \lstDeleteShortInline{|}%
}

\newenvironment{cpplstShortInline}{%
    \lstMakeShortInline[
        style=cppstyle,
        style=lstinlinestyle,
    ]{|}%
}{%
    \lstDeleteShortInline{|}%
}

\newenvironment{llstShortInline}{%
	\lstMakeShortInline[
		style=lstyle,
		style=lstinlinestyle,
	]{|}%
}{
	\lstDeleteShortInline{|}%
}

\newenvironment{pylstShortInline}{%
	\lstMakeShortInline[
	style=pystyle,
	style=lstinlinestyle,
	]{|}%
}{%
	\lstDeleteShortInline{|}%
}

\newenvironment{mlstShortInline}{%
    \lstMakeShortInline[
        style=mstyle,
        style=lstinlinestyle,
    ]{|}%
}{%
    \lstDeleteShortInline{|}%
}

\newenvironment{alstShortInline}{%
    \lstMakeShortInline[
        style=astyle,
        style=lstinlinestyle,
    ]{|}%
}{%
    \lstDeleteShortInline{|}%
}

\newenvironment{amlstShortInline}{%
	\lstMakeShortInline[
	style=amstyle,
	style=lstinlinestyle,
	]{|}%
}{%
	\lstDeleteShortInline{|}%
}

\newenvironment{leglstShortInline}{%
    \lstMakeShortInline[
        style=legstyle,
        style=lstinlinestyle,
    ]{|}%
}{%
    \lstDeleteShortInline{|}%
}

\newenvironment{pinkylstShortInline}{%
	\lstMakeShortInline[
	style=pinkystyle,
	style=lstinlinestyle,
	]{|}%
}{%
	\lstDeleteShortInline{|}%
}

% gebruik voor opmaak van console output
% http://ftp.snt.utwente.nl/pub/software/tex/macros/latex/contrib/fancyvrb/fancyvrb.pdf
% nodig om underline voor input te kunnen gebruiken: \cinput{...}
\usepackage{fancyvrb}
\definecolor{darkgreen}{RGB}{0,100,0}
\newcommand{\cinput}[1]{\color{darkgreen}{\underline{#1}}}
\newenvironment{coutput}{%
	\Verbatim[commandchars=\\\{\}]%
    }{%
    \endVerbatim%
}

%http://ctan.cs.uu.nl/macros/generic/xstring/xstring_doc_en.pdf
\usepackage{xstring}
\makeatletter
\newrobustcmd{\proglink}[1]{%
    \normalexpandarg%
    \StrDel{#1}{\\-}[\bare@link]%
    \href{\progsurl/\bare@link}{\hcode{#1}}%
}
\makeatother
